% PDF 164--204；原书第 135--175 页。
\chapter{大语言模型的适配与使用}
\label{ch:adaptation}
\chaptercredit{4}
\chapterabstract{本章考察按用户需求适配和使用大语言模型的若干方法。首先讨论指令调优，这是让大语言模型适应特定任务和领域的重要步骤。随后重点介绍参数高效微调，这类技术只调整少量参数来改善模型表现，在保持较高准确率的同时节约算力。接着讨论在微调时提高计算资源利用效率的几种策略，说明如何通过量化减轻计算负担，并保持模型的效果。谈完大语言模型的适配，再转向它的使用，介绍终端用户提示大语言模型时涉及的主要概念，并给出提示工程的实际指导。最后的综合教程说明如何在资源受限时微调大语言模型。}

% PDF 164--165 上半页。
\section{引言}
第~\ref{ch:pretraining}~章已经介绍，预训练大语言模型时，可以选用不同的神经网络架构、数据、学习目标或任务。无论预训练方式有何差别，模型通常都能从接触过的语言中学到有用的模式。由于预训练规模空前，这些学到的模式往往使大语言模型具备零样本完成语言任务的能力，表现常常超过专门针对这类任务微调的小型模型（\refcite{ch04-brown}{Brown et al., 2020}）。

许多研究还表明，若针对某些领域或任务进一步适配或微调，大语言模型在相应领域或任务中的能力可以更强。预训练语料中的领域内容往往比较笼统，即使覆盖广泛，通常也只均匀地分配篇幅。由此可见，这类提升并不意外。本章介绍高效适配和使用大语言模型的技术与概念。

\begin{figure}[!htbp]
\centering
% 重绘自原书第 136 页的图 4.1。
\begin{tikzpicture}[
  font=\footnotesize\sffamily,
  box/.style={draw=black!43,fill=yellow!38,align=center,inner sep=2pt,
    minimum width=23mm,minimum height=5.8mm,rounded corners=1.5pt},
  root/.style={box,minimum width=20mm,minimum height=15mm,rounded corners=0pt},
  wire/.style={draw=blue!35!black!65,line width=.55pt},
  flow/.style={wire,-{Stealth[length=1.4mm]}},
  group/.style={draw=blue!35!black!65,densely dashed,rounded corners=4pt,inner sep=3pt}]
\node[root] (adapt) at (0,0) {大语言模型\\适配};
\node[box,fill=orange!20] (it) at (2.92,0) {指令调优};
\node[box,fill=orange!40] (peft) at (5.64,.70) {参数高效微调};
\node[box,fill=orange!40] (ceft) at (5.64,-.70) {计算高效微调};
\node[box] (adapter) at (8.77,1.04) {适配器};
\node[box] (repara) at (8.77,.36) {重参数化};
\node[box] (ptq) at (8.77,-.36) {训练后量化};
\node[box] (qat) at (8.77,-1.04) {量化感知微调};
\draw[wire] (1.20,0) .. controls (1.91,.05) and (1.68,1.34) .. (2.15,1.34);
\draw[wire] (1.20,0) .. controls (1.91,-.05) and (1.68,-1.34) .. (2.15,-1.34);
\draw[wire] (it.east) -- ++(.20,0) |- (peft.west);
\draw[wire] (it.east) -- ++(.20,0) |- (ceft.west);
\draw[wire] (peft.east) -- ++(.32,0) |- (adapter.west);
\draw[wire] (peft.east) -- ++(.32,0) |- (repara.west);
\draw[wire] (ceft.east) -- ++(.32,0) |- (ptq.west);
\draw[wire] (ceft.east) -- ++(.32,0) |- (qat.west);

\node[root,fill=cyan!21] (use) at (0,-3.14) {大语言模型\\使用};
\node[box,fill=blue!24] (chain) at (2.92,-3.14) {提示链};
\draw[wire] (1.20,-3.14) .. controls (1.91,-3.09) and (1.68,-1.98) .. (2.15,-1.98);
\draw[wire] (1.20,-3.14) .. controls (1.91,-3.19) and (1.68,-4.30) .. (2.15,-4.30);
\node[box,fill=blue!8] (zero) at (5.64,-2.99) {零样本提示};
\node[box,fill=blue!8] (few) at (5.64,-3.68) {少样本提示};
\node at (5.64,-2.32) (context) {提示上下文};
\node[group,fit=(zero)(few)(context)] (contextgroup) {};
\node[box,fill=blue!15] (cot) at (8.77,-2.92) {思维链};
\node[box,fill=blue!15] (sc) at (8.77,-3.61) {自一致性};
\node[box,fill=blue!15] (tot) at (8.77,-4.30) {思维树};
\node at (8.77,-2.25) (ensemble) {提示集成与结构};
\node[group,fit=(cot)(sc)(tot)(ensemble)] (ensemblegroup) {};
\draw[flow] (chain.east) -- (contextgroup.west |- chain.east);
\draw[flow] (contextgroup.east |- chain.east) -- (ensemblegroup.west |- chain.east);
\end{tikzpicture}

\caption{本章介绍的概念分类，重点是大语言模型的高效适配和使用。}
\label{fig:ch04-taxonomy}
\end{figure}

为此，本章考察研究文献中有助于理解问题或具有应用前景的工具、技术与流程。图~\ref{fig:ch04-taxonomy} 汇总了本章范围。在大语言模型适配方面，我们估计大多数读者用于微调的经费有限，所以优先介绍参数高效微调和资源高效微调。适应性预训练等需要调节更多参数或耗用更多资源的方法只作简述。基于提示的学习已在第~\ref{ch:prompting}~章专门讨论。在大语言模型的使用方面，本章介绍终端用户进行提示工程时的基本概念，包括提示链和思维链提示。进入微调技术之前，先介绍大语言模型微调常用的指令调优及其基本概念。

% PDF 165--173；原书第 136--144 页。
\section{指令调优}
\label{sec:ch04-instruction}
第~\ref{sec:ch03-prompt-finetuning}~节介绍了大语言模型基于提示的微调。基本流程是使用表示给定任务的有标签问答对数据集，设计提示模板，把问题转换成自然语言指令，再用这些指令微调模型，使其预测人工整理的答案。这个办法很适合完成指定任务。不过，ChatGPT 和 Claude 2 等当时表现先进的聊天机器人，除了擅长训练所针对的问题，还能将能力扩展到微调目标之外的多种主题。怎样调优模型，使它无需示例便能准确完成与训练数据不同的任务？一种方法是指令调优（instruction tuning，IT）。

\subsection{指令调优流程}
从最基本的意义上说，指令调优就是把提示写成模型可执行的自然语言指令，再用这些提示微调语言模型。提示通常包含三个部分，即描述待完成任务的指令、完成任务所需的上下文，以及引出答案的提示。仍以第~\ref{sec:ch03-prompt-finetuning}~节的例子为例。
\begin{quote}
\texttt{Cannot watch this movie. This is [z].}
\end{quote}

做基于提示的微调时，我们收集大量这样的电影评论句子，将“great”或“terrible”分别填入提示，表示正面或负面类别，再用填好的提示调优模型。也可以为分类任务设计另一种模板，给出明确的自然语言指令，使模型照做。这样便无需从完形填空式提示中推断隐含的指示。

\begin{lstlisting}[nolol=true]
## Instruction: Please read the following film review and
determine if it is a positive or negative review. Respond
'positive' or 'negative', according to your classification.
## Review: Cannot watch this movie.
## Response:
\end{lstlisting}

这种表达更明确，也更接近人们习惯的交流方式。在许多指令模板样例上微调后，大语言模型除了能完成电影评论填空，还能学会遵循自然语言指令，并按人们期望的方式生成回答。

这种办法的一个突出特点是，遵循指令的能力可以扩展到微调中没有明确训练过的新任务。指令调优使用以自然语言表述的问题和目标答案，微调大语言模型，覆盖阅读理解、情感分析、翻译、常识推理等多种自然语言任务。任务既可以是有唯一答案的离散任务，例如“\texttt{What is the capital of Canada?}”，也可以是没有唯一正确答案的开放任务，例如“\texttt{Write a poem about pirates:}”。对后者，同一提示可以提供多条示范。

利用指令模板，把这些训练集转换为自然语言指令，调优模型，使它预测与每个训练输入配对的答案。调优以词元为单位进行。先给出完整的提示，通过反向传播调优模型，使它预测回答的第一个词元。随后，给出提示和第一个词元，使模型预测第二个词元，如此继续下去。多项研究（\refcite{ch04-wei-2021}{Wei et al., 2021}；\refcite{ch04-ouyang}{Ouyang et al., 2022}；综述见 \refcite{ch04-zhang-2023c}{Zhang et al., 2023c}）显示，经过指令调优的大语言模型，在微调时未接触过的自然语言处理任务上表现有所改善，可以很好地泛化。

\begin{figure}[!htbp]
\centering
% 重绘自原书第 138 页的图 4.2。
\begin{tikzpicture}[
  font=\footnotesize\sffamily,
  stagebox/.style={draw=black!44,fill=cyan!20,align=center,inner sep=3pt,
    minimum width=31mm,minimum height=12mm,rounded corners=2.4pt},
  result/.style={stagebox,fill=green!23},
  arrow/.style={draw=black!80,-{Stealth[length=1.4mm]},line width=.65pt}]
\node[stagebox,fill=yellow!27,minimum width=24mm] (pre) at (0,0) {预训练\\语言模型};
\node[stagebox] (ptft) at (3.20,1.66) {针对任务 A\\微调};
\node[stagebox,fill=red!19] (prompt) at (3.20,0) {为任务 A 设计零样本\\或少样本提示};
\node[stagebox] (it) at (3.20,-1.66) {在任务 B、C、D、\dots\\上微调};
\foreach \item/\title in {ptft/预训练再微调,prompt/基于提示的学习,it/指令调优}{
  \node[above=1pt of \item] {\title};
  \node[result,right=5mm of \item] (\item-result) {在任务 A 上推理};
  \draw[arrow] (\item.east) -- (\item-result.west);
}
\draw[arrow] (pre.east) to[out=13,in=190] (ptft.west);
\draw[arrow] (pre.east) -- (prompt.west);
\draw[arrow] (pre.east) to[out=-13,in=170] (it.west);
\end{tikzpicture}

\caption{预训练再微调、基于提示的推理和指令调优的总体流程比较。预训练再微调针对单一任务微调，再在同一任务上推理。基于提示的学习设计一条提示，可以包含上下文示例，随后用该提示推理。指令调优在众多数据集和不同任务上微调，使模型的能力可以泛化到训练中未见过的新任务。}
\label{fig:ch04-workflows}
\end{figure}

指令调优与第~3.1.2~节介绍的预训练再微调，以及第~3.2.2~节详细讨论的基于提示的学习和推理密切相关。图~\ref{fig:ch04-workflows} 显示了这些方法之间的关系。

\begin{keypoint}
预训练再微调对大语言模型做监督微调，使其完成特定的机器学习任务，涵盖基于任务头的微调与基于提示的微调。指令调优也做监督微调，但用统一的模板组织来自许多不同自然语言推断任务的输入与输出对。将输入统一写成自然语言提示后，模型学会执行这些任务，也学会正确理解人所写指令的含义。自回归大语言模型原本依据训练数据预测字符串中最可能出现的下一个词元。要让它给出有用、准确的指令回答，这一过程十分重要。
\end{keypoint}

下面讨论研究者如何收集指令调优数据集，并列举文献中使用指令调优实现领域适配的例子。

\subsection{指令调优数据}
指令调优训练集由问答对组成，将不同自然语言处理任务表述成自然语言指令及所需回答。研究者采用过多种方法来构造数据集，主要有以下三类。

\begin{enumerate}
\item 用提示模板把已有的自然语言处理训练数据集转换为指令调优数据。
\item 收集人编写的提示，再由人编写答案。
\item 收集人编写的提示，用已有的大语言模型生成答案。
\end{enumerate}

为了增加训练集的多样性和规模，往往会结合使用这三种方法。下面简述每种方法及其操作方式。图~\ref{fig:ch04-instruction-data} 给出概览。

\subsubsection{用模板转换的数据集做指令调优}
Hugging Face Hub 和 TensorFlow Datasets 等来源提供了大量可用于指令调优的自然语言处理数据集。这些数据集的结构和用途多种多样，一般有一个或多个输入，例如问题、上下文和指令，以及一个或多个可能的目标输出，输出可能按偏好排序。某些任务有明确的是非答案，如情感的正负判断。另一些需要从提供上下文的段落中抽取信息，如开卷问答。还有些任务比较开放，没有唯一正确的答案，如摘要和翻译。任务和主题的多样性，使数据能够涵盖许多相关的自然语言处理任务及领域知识。

\sourcefigure{fig04_03.pdf}{构造指令调优数据集的三种办法。（1）收集不同推理任务的开源数据集，用统一模板将问答转换成指令与输出，再做监督微调。（2）收集大量指令，例如用户发送给 OpenAI API 的查询，请人工标注者编写一个或多个回答，再微调基础模型。（3）收集种子任务及问答对，用大语言模型生成新任务及问答对，筛去质量较低的数据，再用筛选后的问答对微调基础模型。}{fig:ch04-instruction-data}

要用这些数据进行指令调优，可为每个数据集设计模板，将数据从原有结构转成自然语言指令与目标答案。以清单~\ref{lst:ch04-boolq} 中原书取自 GLUE BoolQ 的上下文、问题、答案三元组为例。

\noindent\begin{minipage}{\linewidth}
\begin{lstlisting}[caption={原书的 GLUE BoolQ 样例},label={lst:ch04-boolq}]
passage: "Look What You Made Me Do" is a song recorded by
American singer -songwriter Taylor Swift, released on August
24, 2017 by Big Machine Records as the lead single from her
sixth studio album Reputation (2017). Swift wrote the
song with her producer Jack Antonoff....
question: "did taylor swift write look what you made me do"
label: 1 (yes)
\end{lstlisting}
\end{minipage}\par

数据集的每条记录包含提供上下文的段落、一个关于该段落的问题，以及一个是非答案。可以按图~\ref{fig:ch04-instruction-data} 最上方流程中的模板 1 将它整理成下面的形式。

\noindent\begin{minipage}{\linewidth}
\begin{lstlisting}[caption={问答样例},label={lst:ch04-qa}]
"Look What You Made Me Do" is a song recorded by American singer-
songwriter Taylor Swift, released on August 24, 2017 by Big
Machine Records as the lead single from her sixth studio
album Reputation (2017). Swift wrote the song with her
producer Jack Antonoff... .

Based on the paragraph above answer the following question: did
taylor swift write "look what you made me do ?"
yes
\end{lstlisting}
\end{minipage}\par

训练时，把问题作为初始条件交给模型，再微调模型，使其正确回答“yes”。同一个数据集可以用几个不同的提示模板。这有助于防止模型对单一模板的特定措辞过拟合，也能减轻单一选择可能存在的缺点。

主要通过整理已有数据进行指令调优的代表模型是 FLAN（\refcite{ch04-wei-2021}{Wei et al., 2021}），它由大语言模型 LaMDA（\refcite{ch04-thoppilan}{Thoppilan et al., 2022}）经指令调优得到。研究者从开源数据库收集 62 个有标签的自然语言处理训练数据集，按任务分成 12 组。其中大多数是阅读理解和情感分析等自然语言理解任务，少数是机器翻译和摘要等生成任务。他们为每个数据集设计多个模板。为了检验指令调优模型的泛化能力，每次留出一个任务组，在其他 11 组上调优，再按留出组上的表现评分。与未经微调的 LaMDA 相比，FLAN 的表现大幅提升，尤其是在翻译和闭卷问答任务上。它还优于 GPT-3 等规模大得多、未经指令调优的大语言模型。这表明，要充分发挥模型的能力，指令调优十分重要，无论模型的参数量如何。

\subsubsection{用人工生成的样例做指令调优}
另一种方法是编写或收集一组指令任务，再请人写出答案供模型学习。我们称它们为人工生成的样例。这种方法适合为写诗等自然语言生成任务收集数据。这些任务没有特定的正确答案，但某些输出比其他输出更符合人的偏好。用人工编写的文本训练，有助于模型在回答时采用更接近人的语言表达方式。人工样例的主要优点是，模型设计者可以更细致地选择问题主题和规定回答细节。人工标注耗时且昂贵，因此专门人工制作的数据集，往往比汇集已有数据集得到的数据集更小。

InstructGPT（\refcite{ch04-ouyang}{Ouyang et al., 2022}）是利用人工生成样例的重要模型。研究者收集发送给 OpenAI API 的查询，并补充一些人工编写的问题，以构成输入集。随后雇用一组标注者，为这些查询撰写答案，构成完整数据集。他们采用的训练方法与第~\ref{ch:alignment}~章讨论的基于人类反馈的强化学习密切相关。训练后，标注者往往更喜欢 InstructGPT 的输出。与作为参照的模型比较，InstructGPT 的胜率约为 73\%，FLAN 约为 30\%。这个结果一部分来自训练技术，也反映了微调用的输出由人编写。用人工编写的答案训练的模型，比用多个不同数据集模板化训练的模型更符合人类偏好。

\subsubsection{用大语言模型生成的样例做指令调优}
最后一种指令调优方法，是用大语言模型回答查询，生成数据集。可以用要训练的同一个模型，也可以用另一个具有希望通过微调获得的特性的模型。一般先让模型生成若干答案，再用程序或人工筛选出质量最好的回答。与人工作答相比，这种方法快得多，费用也低得多，但回答所能涵盖的范围受大语言模型能力限制。

如果把答案交回最初生成它们的模型训练，还会遇到明显的数据污染问题。这样能否带来新信息？我们有理由追问，模型用自己生成的回答微调，究竟能学到什么。一种解释是，先让模型为同一个问题生成许多答案，再选出最符合人类判断的一两个，用它们微调模型。这样，模型就更有可能输出人们认为好的答案。即使它不像使用人工生成数据微调那样接触到新信息，也能学到人们希望它表现出的倾向。

为用实验检验相近的想法，\refcite{ch04-wang-2023b}{Wang 等人（2023b）}开发了 Self-Instruct，一种迭代式自举框架。它接收一组种子任务及各任务的示例输入与输出，再借助 GPT-3 根据种子输入生成新的查询。随后把查询交给 GPT-3 作答，将查询与输出配对构造指令调优数据集。研究者使用多种质量筛选方法，确保回答的准确性和多样性。最终结果表明，在自然语言生成任务上，他们的模型明显优于未经调优的 GPT-3，也优于由更小的大语言模型 T5 经指令调优得到的模型。

\subsection{用指令调优做领域适配}
前面主要谈了怎样推理未见过的任务。指令调优也常用于让大语言模型适应特定领域。研究表明，当任务高度专门化，需要理解专业术语或掌握预训练数据以外的信息时，经过领域适配的指令调优模型可以优于通用聊天机器人。这里概述教育、医学和金融领域的例子，这些例子只是指令调优的一部分应用。

在教育领域，\refcite{ch04-zhang-2023d}{Zhang 等人（2023d）}发布了用于写作任务的 Writing-Alpaca。调优后，模型能够对写作提出改进建议，例如纠正语法、使表达更清楚、简化难懂的句子或改写文本。调优数据包含每项任务的一句简短指令、要改正的输入句子以及改正后的输出句子。例如：

\begin{lstlisting}[nolol=true]
###Instruction:
Fix grammatical errors in the text
###Input:
She went to the markt
###Response:
She went to the market
\end{lstlisting}

这样，模型学会发现并改正拼写错误以及其他语法错误。Zhang 等人用约 60\,000 个文本改进样例微调 Llama-7B，涵盖七种文稿编辑任务，并从 Stanford-Alpaca 项目（\refcite{ch04-taori}{Taori et al., 2023}）另外加入 52\,000 条通用指令提示。所得模型显著优于作为参照的基础模型以及用专门程度较低的数据集调优的模型，但仍未达到用数百万个文本编辑样例做预训练再微调的模型的表现。这表明，指令调优可以用规模相对有限的数据得到质量较好的模型，而达到最佳表现有时仍需要更大的数据集。

在医学领域，\refcite{ch04-li-2023b}{Li 等人（2023b）}提出了 ChatDoctor，也是对大语言模型 Llama-7B 进行指令调优的版本。研究者用真实的医患文字聊天记录训练模型，使它能针对患者提出的问题给出行动建议。这个项目用指令提示调优模型，依次完成三项任务。

\begin{enumerate}
\item 用指令提示教模型从患者的问题中提取关键词。
\item 用指令提示引导模型考虑从内部数据库与维基百科取得、与关键词相关的上下文，从中抽取与患者问题有关的信息。
\item 用指令提示让模型考虑抽取出的上下文，并向患者建议接下来可采取的措施。
\end{enumerate}

训练数据由程序从医患对话中抽取，用于构造约 100\,000 个按顺序组织的样例。经过指令调优的模型，在理解患者症状并提供与症状有关的建议方面有明显改善。这也说明，通过串联多条指令，可以完成单条提示难以处理的更复杂任务。医学领域还有经过指令调优的大语言模型，专门用于放射学和肿瘤学等更细分的学科。

金融领域也出现了为不同任务设计的指令调优模型。例如，\refcite{ch04-zhang-2023a}{Zhang 等人（2023a）}开发了 Instruct-FinGPT，自动判断财经标题表达的情感是正面、中性还是负面。

\sourcefigure{fig04_04.pdf}{用模板处理金融数据以调优 Instruct-FinGPT 的例子（\refcite{ch04-zhang-2023a}{Zhang et al., 2023a}）。左侧为句子及对应标签，中间为指令模板，右侧为指令提示及回答。}{fig:ch04-fingpt}

研究者选用文献中标有情感类别的财经标题，拟定十条不同的情感分类指令，再构造带有标注答案的提示。图~\ref{fig:ch04-fingpt} 给出了一些模板转换的例子。这些例子说明了领域适配的用处。许多财经标题使用专门的措辞和句法，需要模型更充分地学习这些特征，才能取得更好的结果。使用这些数据调优 Llama-7B 后，与作为参照的 Llama 模型和 ChatGPT 等当时表现先进的聊天机器人相比，$F_1$ 得分提高超过 20\%。这说明了以指令调优进行领域适配的价值。

以上几例说明了怎样组织数据并用它微调大语言模型，从而改善表现、拓展能力。后续各节将考察高效微调大语言模型的不同办法。

% PDF 173--186；原书第 144--157 页。
\section{参数高效微调}
\label{sec:ch04-peft}
大语言模型的规模对某些有价值的能力的涌现至关重要，却也带来了一些实际限制。最明显的是，怎样在不同环境之间高效传输大型模型，以及怎样存储它们（\refcite{ch04-ding}{Ding et al., 2023}）。如果每种任务、使用场景或应用都要保留多个模型副本，这些困难会进一步加重。即便 BERT 等较小的语言模型也会遇到类似的问题。BERT 是 Google Research 团队在 \refcite{ch04-houlsby}{Houlsby 等人（2019）}的工作中首次提出适配器的 Transformer 模型，下文会详细讨论适配器。

除了部署方面的困难，如果微调策略针对大语言模型的全部参数，模型规模往往会使所需时间和费用超出可承受范围。这样的限制压缩了研究人员与开发人员试验和反复改进调优策略的空间。因此，必须事先投入相当的精力，尽量保证预算所能负担的少数几次试验有效且不出差错，因为出错的代价可能很高。

大语言模型通过预训练获得很强的通用能力，而微调的目标，按定义，是使模型在一个或多个相对狭窄的能力范围内表现得更专门化。如果为了完成目标任务而微调大语言模型的所有参数，部分通用能力也必然发生改变。由此得到的模型，通用性不及只做预训练的模型。因此，需要保留多少版本和副本，可能取决于要覆盖的任务数量。

参数高效微调（parameter-efficient fine-tuning，PEFT）是一组试图解决这些困难的微调技术。一般而言，它希望用尽可能少的可调参数取得尽可能好的任务结果。PEFT 系统通常也考虑保留大语言模型在预训练中得到的有用通用能力。常见做法是针对每项任务调优附加到模型的适配器（adapter）或其他扩展，而大语言模型的预训练参数保持不变（\refcite{ch04-houlsby}{Houlsby et al., 2019}）。

\sourcefigure{fig04_05.pdf}{四种参数高效微调方法的模型架构细节。（a）基于提示的微调方法中的前缀调优。（b）重参数化微调方法中的 LoRA。（c）在 Transformer 的既有网络层中串联适配器。（d）将适配器与 Transformer 的既有网络层并联。}{fig:ch04-peft-types}

依照 \refcite{ch04-hu-2023}{Hu 等人（2023）}的划分，PEFT 方法大体分为四类，如图~\ref{fig:ch04-peft-types} 所示。

\begin{itemize}
\item \textbf{提示方法。}这类方法与基于提示的学习方法有关，用于微调大语言模型。一般目标是构造可学习的提示向量，将其附加到输入嵌入中（提示调优，\refcite{ch04-lester}{Lester et al., 2021}），或附加到预训练语言模型各层的隐藏状态中（前缀调优，\refcite{ch04-li-liang}{Li and Liang, 2021}）。见图~\ref{fig:ch04-peft-types}。
\item \textbf{串联适配器方法。}将额外的神经网络模块与预训练语言模型的既有层串联。以 \refcite{ch04-houlsby}{Houlsby 等人（2019）}为例，适配器网络直接接在 Transformer 的前馈层和注意力层之后。见图~\ref{fig:ch04-peft-types}（c）。
\item \textbf{并联适配器方法。}同样使用小型的附加神经网络模块，但将它们与既有网络的子层并联。在这种配置下，预训练语言模型的中间激活值通过捷径连接传给适配器模块（如 \refcite{ch04-sung}{Sung et al., 2022}）。见图~\ref{fig:ch04-peft-types}（d）。
\item \textbf{重参数化方法。}利用内在维度（intrinsic dimension）的概念，根据给定的微调任务对预训练网络重参数化（\refcite{ch04-aghajanyan}{Aghajanyan et al., 2020}）。LoRA（\refcite{ch04-hu-2021}{Hu et al., 2021}）和 VeRA（\refcite{ch04-kopiczko}{Kopiczko et al., 2023}）属于此类。见图~\ref{fig:ch04-peft-types}（b）。
\end{itemize}

\begin{practicaltip}
选择 PEFT 方法时，主要要权衡计算需求和任务表现。一般来说，计算需求越低，任务表现的损失越大。按前面列出的类别，可以给出一种近似次序，从计算负担较重排到较轻，任务表现也相应降低。

\begin{enumerate}
\item 全参数微调。
\item 基于提示的调优。
\item 串联适配器。
\item 并联适配器。
\item 标准重参数化。
\item 优化重参数化。
\end{enumerate}

标准重参数化指 LoRA 一类方法，采用标准浮点精度和标准优化算法，例如 Adam（\refcite{ch04-kingma}{Kingma and Ba, 2017}）。优化重参数化指 QLoRA 一类方法，为计算开销优化了浮点精度，并用优化算法有选择地更新参数。
\end{practicaltip}

下面考察几种常用 PEFT 方法在架构和任务表现上的得失，包括重参数化、串联适配器和并联适配器。基于提示的学习类 PEFT 方法，其基础已在第~\ref{ch:prompting}~章详细讨论。读者可查阅 Hugging Face 整理的 PEFT 方法\footnote{\url{https://github.com/huggingface/peft}}，从实际使用的角度了解它们。本章还提供带代码的教程，将比较其中几种方法的收益和开销。

\subsection{适配器}
\label{sec:ch04-adapters}
适配器是在待适配的预训练模型各层中加入的小型神经网络模块。所谓“小型”，是相对于语言模型的参数数量而言（\refcite{ch04-houlsby}{Houlsby et al., 2019}）。从实际使用来看，适配器有以下优点。

\begin{itemize}
\item \textbf{任务表现。}在许多任务上，适配器只需调节少得多的参数，就能达到接近全参数微调的表现。
\item \textbf{按任务独立调优。}适配器是纳入 Transformer 模型各层的任务专用模块，可以依次开发。要让语言模型适应多个任务，可以逐个任务调优，无需同时在代表所有任务的数据上训练。开发人员可以针对各任务的具体目标优化，避免只依赖所有任务上的综合指标。这一特点能减轻对原预训练语言模型做多任务全量或部分微调所带来的灾难性遗忘（\refcite{ch04-pfeiffer}{Pfeiffer et al., 2021}）。
\item \textbf{可扩展的部署。}适配器的参数数目通常只占目标语言模型的一小部分，因此可在常规计算基础设施上部署任务专用的适配器模块。
\end{itemize}

将适配器与传统微调和基于特征的迁移调优作比较，有助于理解它的做法。考虑参数为 $w$ 的神经网络 $\phi_w(x)$。传统微调针对每项任务调整原参数 $w$，因此每项任务都需要保存一份新的 $w$，增加了存储需求。基于特征的迁移则用新的函数 $X_v$ 重新组织模型函数，得到 $x_v(\phi_w(x))$，其中只调优新的任务专用参数 $v$。原模型参数 $w$ 保持不变，因此存储开销较小。

尽管计算开销远低于全参数微调，适配器已能取得相当甚至更好的表现。这归功于它组合任务专用部分 $x_v$ 和原模型参数 $w$ 的方式。设定新任务专用参数的初值 $v_0$，使新模型函数尽量接近原函数，即 $\psi_{w,v_0}(x)\approx\phi_w(x)$，训练时只调优 $v$（\refcite{ch04-houlsby}{Houlsby et al., 2019}）。

微调大语言模型时，通常有 $|v|\ll|w|$。也就是说，适配器中的可调参数数量只占原大语言模型参数数量的一小部分。对 \refcite{ch04-houlsby}{Houlsby 等人（2019）}提出的适配器架构，可训练参数数目为 $2md+d+m$。其中，$d$ 是 Transformer 层前馈投影输出的原始特征维数，$m$ 是为适配器层选择的瓶颈维数。选用较小的 $m$，即可减少任务微调所需的附加参数。\refcite{ch04-houlsby}{Houlsby 等人（2019）}报告，在实际微调中，即便参数数目只有原预训练模型的 0.5\%，也能取得良好的效果。

\subsubsection{串联适配器}
串联适配器与预训练网络的既有层串联。这种 PEFT 方法将计算改写为
\begin{equation}
H_o\gets H_o+f(H_oW_{\mathrm{down}})W_{\mathrm{up}}.
\label{eq:ch04-series}
\end{equation}
这里，$H_o$ 是给定网络层的输出。装入串联适配器后，输出先经 $W_{\mathrm{down}}\in\mathbb{R}^{d\times r}$ 向下投影到较低维度，其中 $r$ 为适配器瓶颈大小，通常较小。对降维投影应用非线性函数 $f$，再用 $W_{\mathrm{up}}\in\mathbb{R}^{d\times r}$ 将输出向上投影，恢复到 $H_o$ 的原始维数。%
\footnote{译注：原文将 $W_{\mathrm{up}}$ 的形状也印为 $d\times r$。按式~\eqref{eq:ch04-series} 的右乘顺序和恢复到 $d$ 维的要求，$W_{\mathrm{up}}$ 的形状应为 $r\times d$。}
$W_{\mathrm{down}}$、$f$ 和 $W_{\mathrm{up}}$ 这三部分组成串联适配器，在适配器调优时微调。

图~\ref{fig:ch04-series} 给出适配器的放置位置。一处紧接在多头注意力层之后的前馈层后面，另一处紧接在 Transformer 输出归一化之前的两个前馈层后面。\refcite{ch04-hu-2023}{Hu 等人（2023）}指出，对于某些任务，这些位置未必是最佳选择。他们评估了适配器位置对任务表现的影响。在数学推理任务上，只把适配器放在前馈层之后，比只放在多头注意力层后，或在多头注意力层与前馈层后都放，表现更好。这与 \refcite{ch04-pfeiffer}{Pfeiffer 等人（2021）}提出的更高效的适配器变体一致。

\sourcefigure{fig04_06.pdf}{Transformer 层中适配器位置的架构图。适配器装在两处。第一处位于紧接多头注意力层的前馈投影之后，第二处位于两个前馈层之后。适配器的主要特点是瓶颈结构。输入先被投影到更小的特征空间，经过非线性变换后，再投影回原来的输入维数。}{fig:ch04-series}

除了放置位置，初次向下投影时采用的瓶颈大小 $r$，也是设计适配器时重要的超参数。通常，$r$ 太小会限制从适配器输入层到瓶颈层时能够保留的有用信息。$r$ 太大可能改善任务表现，却会降低微调的参数效率。不过，\refcite{ch04-hu-2023}{Hu 等人（2023）}发现，$r$ 过大也可能使任务表现变差。

串联适配器通常能有效减少微调时的计算量。但由于预测时输入必须经过额外的串联层，有研究报告说，这会给推理带来不可忽略的开销。

\subsubsection{并联适配器}
最初提出适配器并联，是为解决多语言机器翻译的性能下降问题（\refcite{ch04-zhu}{Zhu et al., 2021}）。Zhu 等人希望借助并联适配器缩小表现差距。当时，多个分别处理两种语言的机器翻译模型，表现优于一个多语言机器翻译模型。他们在三个多语言机器翻译基准数据集中，有两个验证了这一办法的效果。图~\ref{fig:ch04-parallel} 给出 Zhu 等人所用的并联适配器架构和放置位置。

并联适配器将计算改写为
\begin{equation}
H_o\gets H_o+f(H_iW_{\mathrm{down}})W_{\mathrm{up}},
\label{eq:ch04-parallel}
\end{equation}
其中，$H_i$ 和 $H_o$ 分别是相应网络层及适配器的输入与输出。

与串联适配器相比，将适配器与主干网络并联的一项主要优点是，训练的计算负担可以大幅降低。可调参数空间已经显著缩小，并且在更新参数时，通常无需沿预训练语言模型反向传播以求得梯度（\refcite{ch04-sung}{Sung et al., 2022}）。

\sourcefigure{fig04_07.pdf}{为提高多语言机器翻译表现而微调的并联适配器，其位置与架构如图所示。架构中不含适配器的 Transformer 部分作为多语言模型进行预训练，同时用双语语料库微调各层适配器，以提高相应语言对的机器翻译表现。}{fig:ch04-parallel}

\subsection{重参数化}
\label{sec:ch04-reparameterization}
与适配器不同，重参数化方法通常无需加入额外的神经网络模块，因此可以减轻适配器带来的推理延迟（\refcite{ch04-hu-2021}{Hu et al., 2021}）。这些方法所利用的基本事实是，许多自然语言处理任务的内在维度比预训练模型的维度低几个数量级，因此可用相对较少的参数有效地针对多种任务微调。

重参数化方法一般先做某种秩分解，再在学习阶段优化预训练模型中较高维表示的低秩表示。下面介绍三种用于大语言模型参数高效微调的代表方法，即低秩适配器（Low-Rank Adapters，LoRA）（\refcite{ch04-hu-2021}{Hu et al., 2021}）、克罗内克适配器（Kronecker Adapters，KronA）（\refcite{ch04-edalati}{Edalati et al., 2022}），以及基于向量的随机矩阵适配（Vector-based Random Matrix Adaptation，VeRA）（\refcite{ch04-kopiczko}{Kopiczko et al., 2023}）。

\subsubsection{低秩适配器}
\label{sec:ch04-lora}
低秩适配器微调通过学习低秩矩阵，近似表示针对当前任务所需的参数更新（\refcite{ch04-he-2022}{He et al., 2022}）。\refcite{ch04-he-2022}{He 等人（2022）}指出，LoRA 有四项主要优点。

\begin{itemize}
\item 同一个预训练模型可以由许多自然语言处理任务共享，每个任务分别学习自己的 LoRA 模块。替换学到的低秩矩阵即可切换任务，大幅减少存储开销和任务切换开销。
\item 优化仅涉及加入的低秩矩阵，无须调整预训练语言模型的全部参数，因此训练所需算力和硬件资源可减至原来的约三分之一。
\item LoRA 微调采用线性结构，因此学到的低秩矩阵可以与预训练语言模型的固定权重合并，不会增加推理延迟。
\item LoRA 旨在寻找所微调任务的低维表示，因此可以与不针对秩优化的调优方法分别使用，也可以与其中许多方法组合使用。
\end{itemize}

这些优点如何实现？LoRA 针对每项需要微调的自然语言处理任务，优化一个小得多的参数集 $\Theta$。考虑全参数微调所优化的下列目标。
\begin{equation}
\max_{\Phi}\sum_{(x,y)\in Z}\sum_{t=1}^{|y|}
  \log\bigl(P_{\Phi}(y_t\mid x,y_{<t})\bigr).
\label{eq:ch04-full-tuning}
\end{equation}
这里，$Z=\{(x_i,y_i)\}_{i=1,\ldots,N}$ 是给定自然语言处理任务的 $N$ 个上下文—目标对。在摘要任务中，$x_i$ 是待摘要的全文，$y_i$ 是相应的摘要。微调时，以预训练模型权重初始化 $\Phi_0$，沿梯度反复更新到 $\Phi_0+\Delta\Phi$，以最大化式~\eqref{eq:ch04-full-tuning}。

全参数微调直接更新预训练模型权重，因此如上所述，实际部署多个任务模型时可能因开销过高而难以实施。\refcite{ch04-hu-2021}{Hu 等人（2021）}提出用 $\Delta\Phi=\Delta\Phi(\Theta)$ 估计特定于任务的参数更新 $\Delta\Phi$。由于自然语言处理任务的内在维度相对于预训练模型较低，可以取 $|\Theta|\ll|\Phi_0|$。现在可以通过对 $\Theta$ 优化下式来估计 $\Delta\Phi$。
\begin{equation}
\max_{\Theta}\sum_{(x,y)\in Z}\sum_{t=1}^{|y|}
  \log\bigl(p_{\Phi_0+\Delta\Phi(\Theta)}(y_t\mid x,y_{<t})\bigr).
\label{eq:ch04-lora-objective}
\end{equation}

从算法角度看，\refcite{ch04-hu-2021}{Hu 等人（2021）}关注 Transformer 架构中的全连接层。他们假设，可以把预训练权重矩阵 $W_0\in\mathbb{R}^{d\times k}$ 的更新限制为较低秩的分解，即 $W_0+\Delta W=W_0+BA$，其中 $B\in\mathbb{R}^{d\times r}$、$A\in\mathbb{R}^{r\times k}$，且 $r$ 远小于 $k$ 和 $d$。$k$ 和 $d$ 是满秩的全连接层权重矩阵的维数。训练时只有 $A$ 和 $B$ 的参数可学习，$W_0$ 冻结。因此，对于给定输入 $x$，前向传播的输出为
\begin{equation}
h=W_0x+\Delta W x=W_0x+BAx.
\label{eq:ch04-lora-forward}
\end{equation}
原书说，$A$ 和 $B$ 随机初始化（LoRA 原论文采用高斯随机初始化），同时满足 $\Delta W=BA=0$。%
\footnote{译注：Hu 等人的 LoRA 论文第 4.1 节明确给出 $A$ 用高斯随机值初始化，$B$ 初始化为零。见 \url{https://arxiv.org/html/2106.09685}。}
初始化后，如前所述，通过优化式~\eqref{eq:ch04-lora-objective} 来近似 $\Delta W$。

理论上，LoRA 可以用于任何全连接层的权重矩阵。\refcite{ch04-hu-2021}{Hu 等人（2021）}在最初的应用中只适配自注意力权重。原书这里列出键、值权重矩阵 $W_k$ 和 $W_v$。%
\footnote{译注：LoRA 原论文第 4.2 节举出的部署例子适配的是查询、值投影矩阵 $W_q$ 和 $W_v$，第 7.1 节另比较了不同注意力矩阵的选择。}
尽管使用范围有限，将 LoRA 用于 GPT-3（\refcite{ch04-brown}{Brown et al., 2020}）后，在三个常用基准 GLUE、WikiSQL 和 SAMSum 上，表现达到或超过全参数微调。

\refcite{ch04-hu-2021}{Hu 等人（2021）}还通过实验说明，用较低的秩 $r$ 适配多种层的矩阵，比只适配少数几种层的矩阵而采用较高秩，在参数数量和任务质量之间取得的折中更好。这表明以任务的内在维度解释微调的基本假定，在实验条件下成立（\refcite{ch04-hu-2021}{Hu et al., 2021}）。

LoRA 的成效迅速催生了一些后续研究和应用。例如，LoRAHub（\refcite{ch04-huang}{Huang et al., 2023}）着力改善 LoRA 适配器模块在实际应用中组合使用的能力。第~\ref{sec:ch04-quantization}~节还会看到，将量化方法引入 LoRA，也是逐步降低大语言模型微调和部署开销的一项进展（\refcite{ch04-dettmers-2023}{Dettmers et al., 2023}）。

LoRA 可以用于多种自然语言处理任务，即使模型规模很大，也能以较高的计算和数据效率完成这些任务，因此仍是一种常用的微调方法。后续研究在 LoRA 的低秩分解思想上作了改进，希望进一步提高参数效率与任务表现。例如，AdaLoRA（\refcite{ch04-zhang-2023b}{Zhang et al., 2023b}）按一项基于差异的重要性指标，针对给定任务动态分配总参数预算，有选择地更新微调参数。QLoRA（\refcite{ch04-dettmers-2023}{Dettmers et al., 2023}）则通过降低预训练语言模型的浮点精度实现量化，进一步提高梯度反向传播时的计算效率，详见第~\ref{sec:ch04-quantization}~节。

这些低秩方法很有前景，但在某些场景中，秩太低不足以捕捉任务所需的重要模式，LoRA 的秩分解还有改进余地。下面讨论一种参数效率与 LoRA 相近、却无需 LoRA 的低秩假定的方法，即 KronA（\refcite{ch04-edalati}{Edalati et al., 2022}）。

\subsubsection{克罗内克适配器}
克罗内克适配器由 \refcite{ch04-edalati}{Edalati 等人（2022）}提出，利用克罗内克积（Kronecker product）分解进行参数高效微调，避免对自然语言处理任务内在维度作较强的假定。此前已有方法利用克罗内克积（\refcite{ch04-edalati}{Edalati et al., 2022}），例如 Compactor（\refcite{ch04-mahabadi}{Mahabadi et al., 2021}），它将秩为 1 的矩阵和克罗内克积结合，提高微调的参数效率。这些方法的任务表现良好，但训练与推理的计算效率都低于 KronA（\refcite{ch04-edalati}{Edalati et al., 2022}）。KronA 通过优化有关计算来改进这类重新分解方法（见图~\ref{fig:ch04-krona}）。一般而言，矩阵 $A$ 和 $B$ 的克罗内克积为
\begin{equation}
W=A\otimes B=
\begin{bmatrix}
a_{11}B & \cdots & a_{1n}B\\
\vdots & \ddots & \vdots\\
a_{m1}B & \cdots & a_{mn}B
\end{bmatrix},
\label{eq:ch04-kronecker}
\end{equation}
其中，$W$ 是 $A$ 和 $B$ 的克罗内克积所得的分块矩阵，$m$、$n$ 分别是 $A$ 的行数和列数。\refcite{ch04-edalati}{Edalati 等人（2022）}采用更高效的计算，无需直接构造 $W$。
\begin{equation}
(A\otimes B)x=\gamma\bigl(B\eta_{b_2\times a_2}(x)A^{\mathsf T}\bigr).
\label{eq:ch04-kronecker-product}
\end{equation}
这里，$(A\otimes B)x$ 是矩阵 $A$ 和 $B$ 的克罗内克积与输入向量 $x\in\mathbb{R}^{d_h}$ 相乘的结果。$d_h$ 为输入嵌入维数，$A^{\mathsf T}$ 为矩阵 $A$ 的转置。$\eta_{m\times n}(x)$ 将向量 $x$ 变成 $m\times n$ 的矩阵，$\gamma(x)$ 则把矩阵的各列依次叠放成向量。

在图~\ref{fig:ch04-krona}（a）中，对于给定的输入 $X$，输出 $Y$ 为
\begin{equation}
Y=XW+sX[A_k\otimes B_k],
\label{eq:ch04-krona-output}
\end{equation}
其中，$s$ 为缩放因子，$W$ 是预训练语言模型被冻结的权重。因此，对给定自然语言处理任务使用 KronA 微调后，权重为
\begin{equation}
W_{\mathrm{tuned}}=W+s[A_k\otimes B_k].
\label{eq:ch04-krona-weights}
\end{equation}

\sourcefigure{fig04_08.pdf}{（a）KronA 与（b）$\mathrm{KronA}^{\mathrm B}_{\mathrm{res}}$ 的架构。$\otimes$ 表示矩阵 $A$ 与 $B$ 的克罗内克积。$\mathrm{KronA}^{\mathrm B}_{\mathrm{res}}$ 含有残差连接 $S_{\mathrm{res}}$，去掉它便得到 $\mathrm{KronA}^{\mathrm B}$。}{fig:ch04-krona}

$A_k$ 和 $B_k$ 替代了 LoRA 架构中的向下投影和向上投影（见图~\ref{fig:ch04-peft-types}（b））。和 LoRA 一样，它们可与大语言模型的权重合并。冻结大语言模型权重，并在最后合并权重，可实现高效微调而不增加推理延迟。

\refcite{ch04-edalati}{Edalati 等人（2022）}还提出一种与预训练语言模型前馈网络模块并联的 KronA 适配器架构，记为 $\mathrm{KronA}^{\mathrm B}$，以及在同一架构中加入残差缩放因子以进一步提高任务表现的版本。不过，这两种架构在微调和推理时的计算效率都较低，本书不再详细讨论。感兴趣的读者可参阅 \refcite{ch04-edalati}{Edalati 等人（2022）}。

与其他 PEFT 方法相比，KronA 在任务表现和计算方面如何？\refcite{ch04-edalati}{Edalati 等人（2022）}报告，将 KronA 用于 T5（\refcite{ch04-raffel}{Raffel et al., 2020}）并在 GLUE 基准上评估，其平均表现优于全参数微调、Compactor（\refcite{ch04-mahabadi}{Mahabadi et al., 2021}）、BitFit（\refcite{ch04-zaken}{Zaken et al., 2022}）、LoRA（\refcite{ch04-hu-2021}{Hu et al., 2021}）及 \refcite{ch04-he-2022}{He 等人（2022）}提出的并联适配器。与全参数微调相比，KronA 取得上述表现时，训练时间缩短了 25\%（LoRA 为 28\%），也没有增加推理延迟。在同一基准上，$\mathrm{KronA}^{\mathrm B}$ 和 $\mathrm{KronA}^{\mathrm B}_{\mathrm{res}}$ 均优于 KronA。

\subsubsection{基于向量的随机矩阵适配}
LoRA 一类的重参数化方法能将可训练参数数目降至原来的万分之一，GPU 显存需求最多可降至原来的约三分之一。但有些使用场景除了需要针对任务适配大语言模型，可能还需在各任务上适配每位用户，例如个性化助理、个性化推荐和边缘设备。\refcite{ch04-kopiczko}{Kopiczko 等人（2023）}指出，即使采用 LoRA 达到的参数效率，在生产环境运行时，存储和网络开销仍可能难以承受。受到 \refcite{ch04-aghajanyan}{Aghajanyan 等人（2020）}研究自然语言处理任务微调中内在维度的启发，他们提出了基于向量的随机矩阵适配，即 VeRA（图~\ref{fig:ch04-vera}）。与 LoRA 相比，VeRA 可进一步将微调时的可调参数数目降至约十分之一（图~\ref{fig:ch04-vera-comparison}），因而能减轻规模日益增大的语言模型在实际使用时所遇到的困难。

\sourcefigure{fig04_09.pdf}{LoRA 与 VeRA 的理论显存需求（字节）和可训练参数数目的比较，秩取 $1$、$16$、$256$，模型为 RoBERTa-base、RoBERTa-large 和 GPT-3。原图假设两种方法都只适用于 Transformer 自注意力头的查询层和键层。VeRA 所需显存及可训练参数数目均较低。显存字节数和可训练参数数目用以 10 为底的对数刻度绘制。VeRA 的参数数目为 $|\Theta|=L_{\mathrm{tuned}}(d_{\mathrm{model}}+r)$，LoRA 为 $|\Theta|=2L_{\mathrm{tuned}}d_{\mathrm{model}}r$。其中 $L_{\mathrm{tuned}}$、$d_{\mathrm{model}}$ 和 $r$ 分别是微调的层数、这些层的维数和适配器矩阵的秩。}{fig:ch04-vera-comparison}

这种方法与 LoRA（图~\ref{fig:ch04-peft-types}（b））一样使用一对随机初始化的矩阵 $A$ 和 $B$，初始化细节见下文。微调时，这两个矩阵冻结并由 Transformer 的所有层共享。为了学习微调所得的权重更新 $\Delta W$，VeRA 使用一对可调的缩放向量，即图~\ref{fig:ch04-vera} 中的 $d$ 和 $b$，使冻结的权重矩阵适应给定的自然语言处理任务。这一设计只存储较小的任务专用向量模块，省去保存 LoRA 重参数化矩阵的开销，因此可以在同一个计算节点上存放更多适配后的大语言模型版本。

\sourcefigure{fig04_10.pdf}{VeRA 适配器各组件的架构。与图~\ref{fig:ch04-peft-types}（b）不同，VeRA 冻结 $A$ 和 $B$，并由网络各层共享。微调时只训练向量 $d$ 和 $b$，大幅减少可调参数数量。}{fig:ch04-vera}

回到式~\eqref{eq:ch04-lora-forward}，其中通过调优两个低秩分解矩阵 $A$ 和 $B$ 的乘积来得到 $\Delta W$。VeRA 对给定输入 $x$ 的计算为
\begin{equation}
h=W_0x+\Delta W x=W_0x+\Lambda_b B\Lambda_d A x.
\label{eq:ch04-vera-forward}
\end{equation}
与式~\eqref{eq:ch04-lora-forward} 不同，$A$ 和 $B$ 随机初始化后被冻结，由 Transformer 各层共享。VeRA 不要求 $A$ 和 $B$ 相对大语言模型一定是低秩的，不过这些矩阵的秩增大时，可训练参数数目线性增加。如下文所见，VeRA 的参数数目比 LoRA 少一个数量级以上，同时仍有良好的任务表现。缩放向量 $b$ 和 $d$ 分别写成对角矩阵 $\Lambda_b$ 和 $\Lambda_d$，微调时可训练。$b$ 初始化为全零向量，$d$ 的各元素都初始化为同一非零值。它们在不同层中分别调整矩阵 $A$ 和 $B$ 的行、列大小，使模型适应所关注的自然语言处理任务。

如前所述，VeRA 中的矩阵 $A$ 和 $B$ 是随机初始化的，只需记录重现这些矩阵的随机数生成器种子。因此，VeRA 的存储和内存需求限于这个随机种子与训练得到的向量 $b$ 和 $d$，比 LoRA 少得多，见图~\ref{fig:ch04-vera-comparison}。VeRA 用 Kaiming 初始化（\refcite{ch04-he-2015}{He et al., 2015}）初始化矩阵，使矩阵的方差在不同秩下保持一致，从而减少了为不同的秩分别调节学习率的需要，也节省了训练时间。

考虑到 VeRA 的参数数目比 LoRA 低一个数量级，它与其他 PEFT 方法相比仍有良好的任务表现。在 GLUE 基准上，使用 RoBERTa-base 时，VeRA 仅略逊于 LoRA，使用 RoBERTa-large 时，两者表现相当。在 E2E 基准上比较两者时，使用 GPT-2 的 VeRA 在五项 E2E 任务中的四项上优于 LoRA。

接下来考察提高大语言模型适配和微调效率的其他方法。它们着眼于改进数据表示或使用更高效的优化函数，以减少训练所需内存。许多方法还可以与 PEFT 结合使用，进一步提高效率。

% PDF 186--191；原书第 157--162 页。
\section{计算高效微调}
\label{sec:ch04-compute-efficient}
PEFT 只训练模型中的一小部分参数，以降低微调大语言模型的开销。计算高效微调（compute-efficient fine-tuning）着眼于量化方法，减少在给定参数数量下微调或推理所需的内存。这些方法通常能在训练和推理开销与任务表现之间取得更好的折中。与 LoRA 等常用方法相比，有些方法会损失一些任务表现，另一些则既提高计算资源的利用效率，又改善任务表现，达到或接近当时的最佳水平。

\subsection{大语言模型量化}
\label{sec:ch04-quantization}
量化（quantization）是一种模型压缩技术，用信息位数较少的形式表示参数，从而缩小模型的总体积（\refcite{ch04-zhao}{Zhao et al., 2023}）。这会减少推理所需的计算资源。量化通常施加在 Transformer 注意力层和前馈层的参数权重上，因为这些层的矩阵乘法运算占大语言模型推理内存消耗的 95\% 以上。改变其中使用的数据类型，可以大幅减少内存消耗（\refcite{ch04-dettmers-2022}{Dettmers et al., 2022}）。

\begin{table}[!htbp]
\centering
\caption{大语言模型中的常用数据类型。表中区分源自其他计算领域的标准数据类型和为机器学习优化的表示，并列出其他常用名称和存储所需位数。}
\label{tab:ch04-types}
\small
\begin{tabularx}{\linewidth}{@{}l l X r@{}}
\toprule
数据类型 & 类型来源 & 其他名称 & 位数\\
\midrule
float32 & 标准 & FP32，单精度浮点格式 & 32\\
float16 & 标准 & FP16，半精度浮点格式 & 16\\
bfloat16 & 机器学习优化 & BF16，Brain 浮点格式 & 16\\
INT8 & 标准 & — & 8\\
INT4 & 标准 & — & 4\\
NF4 & 机器学习优化 & — & 4\\
\bottomrule
\end{tabularx}
\end{table}

数据精度与计算效率之间需要权衡。表~\ref{tab:ch04-types} 列出神经网络，尤其是大语言模型常用的几种数据类型的位宽。例如，将参数权重从 float32 量化为 int8，可以将所需内存降至接近四分之一，实际幅度取决于哪些层和参数不参与量化。这个差别很重要，因为一些模型推理时需要的工作内存，远超多数当时先进 GPU 所能提供的显存。例如，具有 1750 亿参数的 GPT-3 在 float16 精度下需要 325~GB 的存储空间，只能在复杂的多 GPU 集群上运行，无法在更常见的硬件上运行，例如单张配备 80~GB 显存的 NVIDIA A100（\refcite{ch04-frantar}{Frantar et al., 2023}）。

大语言模型的量化大体有两种方案，即训练后量化（post-training quantization，PTQ）与量化感知训练（quantization-aware training，QAT）。先考察大语言模型上几项有影响的训练后量化应用，重点是能够把推理的资源消耗降到 NVIDIA A100 或 NVIDIA A600 等常见硬件可承受的范围，并能取得接近同一模型未量化版本的任务表现的方法。

随后讨论微调中的量化感知训练。这时，对微调适配器做量化感知训练，预训练大语言模型本身不参与这种训练。这类应用也提高了大语言模型推理的计算资源利用效率，使预算有限或受到其他资源约束的实践者能够使用它们，例如在微控制器或边缘计算场景中。

\subsubsection{训练后量化}
\label{sec:ch04-ptq}
顾名思义，训练后量化在预训练之后施加于大语言模型。其目标通常是降低推理的内存需求，同时保持原模型的任务表现。朴素量化对权重几乎不加区分地使用低精度数据类型，已被证明对较小的语言模型有效。但大语言模型的参数量超过 67 亿时，观测到任务表现骤降，见图~\ref{fig:ch04-outlier}（\refcite{ch04-dettmers-2022}{Dettmers et al., 2022}）。这一现象与离群特征（outlier feature）的出现有关。它们表现为网络隐藏激活中的较大数值。\refcite{ch04-dettmers-2022}{Dettmers 等人（2022）}最早在大语言模型量化的背景下描述了这一现象。

要保留这些重要离群特征的表示精度，同时满足推理预算，是一个难点。\refcite{ch04-dettmers-2022}{Dettmers 等人（2022）}提出 \texttt{LLM.int8()}，对目标特征中的 99.9\% 逐向量做 INT8 量化，将离群特征单独分离，在矩阵乘法中保持 16 位精度。这增加了量化的复杂性，但这种有选择的混合精度方案使 BLOOM-176B 的推理内存需求减半，并能有效保留原大语言模型的任务表现，如图~\ref{fig:ch04-outlier} 所列的几项基准任务所示。

SqueezeLLM 采用另一种方法，以特征的二阶信息搜索最合适的位精度，从而保留离群特征和其他对精度变化敏感的特征。该方法逐层施加量化，精度可低至 3 位。与原大语言模型相比，推理速度最高可达 2.3 倍，任务表现损失很小（\refcite{ch04-kim}{Kim et al., 2023}）。

ZeroQuant 采用粒度更细的量化方式，分别为权重与激活值采用不同的量化方案，并使用一种知识蒸馏机制抵消任务表现的下降。它在准确率损失很小的情况下，最高可将推理效率提高到原来的 5 倍（\refcite{ch04-yao-2022}{Yao et al., 2022}）。

除了上述方法，GPTQ 也是一种常用的训练后量化方案。和前面的方法一样，GPTQ 用特征的二阶信息搜索量化的最合适位宽。由于有选择地量化权重，支持 4 位、3 位和 2 位的极低位宽，GPTQ 可以让具有 1760 亿参数的 BLOOM 模型在单张 NVIDIA A100 上运行，推理效率最高可达原来的 4.5 倍。\refcite{ch04-liu}{Liu 等人（2023）}则用知识蒸馏技术，尝试改善 3 位精度下的量化效果。

\sourcefigure{fig04_11.pdf}{三种数值存储精度或量化方案用于不同参数规模的语言模型，在多项自然语言推断任务中的表现。参数量从 27 亿增加到 67 亿时，表现大幅下降，这归因于离群特征的出现。\texttt{LLM.int8()} 对大语言模型权重做 8 位量化，同时保留这些特征，因此，即便存储参数权重所用精度减半，其表现仍接近完整的 16 位语言模型。}{fig:ch04-outlier}

\subsubsection{量化感知训练}
\label{sec:ch04-qat}
如第~\ref{sec:ch04-ptq}~节所述，训练后量化在学习过程中并不明确最小化量化操作引入的损失。不过，\texttt{LLM.int8()}、GPTQ 和 SqueezeLLM 等方法以较细的粒度施加训练后量化，在一定程度上也利用了学习过程中的信息。大语言模型预训练采用训练后量化的一个主要原因，是避免训练循环每次迭代时量化大量参数，显著增加计算开销。因此，许多研究试图结合量化带来的推理效率和 PEFT 带来的训练效率。与全参数微调相比，这能减少在训练时加入量化所产生的计算开销。

正如训练后量化使大语言模型能在更易获得的硬件上推理，量化感知训练也可减少微调开销，使之能在更易获得的硬件上进行（\refcite{ch04-dettmers-2023}{Dettmers et al., 2023}）。下面重点介绍原书认为有前景的基于 PEFT 的量化感知训练方法，考察它们能在多大程度上减少微调开销，以及与未经量化的 PEFT 相比，能在多大程度上保留任务表现。%
\footnote{译注：原书这里说要介绍三种方法，本小节实际列出 QLoRA 与 LoftQ 两种。}

\paragraph*{QLoRA}
\refcite{ch04-wortsman}{Wortsman 等人（2023）}提出了把推理时量化所带来的一些效率收益引入训练的思路和建议。在此基础上，QLoRA（\refcite{ch04-dettmers-2023}{Dettmers et al., 2023}）成为大语言模型中采用较广的量化感知训练方法之一。概括地说，QLoRA 先对给定大语言模型采用一种新的 4 位量化，然后在微调中冻结其参数。这项工作引入了一种新数据类型 NF4，即 4 位 NormalFloat。对于服从正态分布的张量数据，NF4 的量化精度优于 4 位整数或 4 位浮点数。量化后，LoRA 权重更新的梯度通过冻结的 4 位量化大语言模型反向传播，因此量化引入的误差会纳入微调过程。

QLoRA 除了采用上述 NF4 数据类型，还使用二次量化（double quantization）进一步减少量化常数引入的内存开销，并采用分页优化器（paged optimizer）。这些办法大幅提高了微调时的计算效率。\refcite{ch04-dettmers-2023}{Dettmers 等人（2023）}结合这三项方法，使用 LoRA 和 Alpaca、FLAN v2 数据集，对具有 650 亿参数的 Llama 大语言模型进行指令微调。他们只在一张 GPU 上微调 24 小时，所得模型即达到 ChatGPT 在相应评估中表现的 99.3\%。实际微调时，16 位精度的全参数微调需要超过 780~GB 的 GPU 显存，QLoRA 所需显存不足 48~GB，同时仍可保持接近当时最佳水平的任务表现。

\paragraph*{LoftQ}
\refcite{ch04-li-2023a}{Li 等人（2023a）}指出，预训练语言模型的量化会影响 LoRA 的微调结果，在极低位宽下尤为明显。为了减轻低位宽量化造成的精度偏差，他们提出量化感知训练技术 LoftQ，试图缩小大语言模型的原始权重矩阵与量化、低秩权重近似共同得到的权重矩阵之间的差距。

这一优化问题表述为 Frobenius 范数的最小化。
\begin{equation}
\min_{Q,A,B}\lVert W-Q-AB^{\mathsf T}\rVert_{f'},
\label{eq:ch04-loftq}
\end{equation}
其中，$\lVert\cdot\rVert_{f'}$ 表示 Frobenius 范数，$W$ 表示原始参数权重，$Q$ 表示量化后的权重，$AB^{\mathsf T}$ 表示 LoRA 所用的低秩项。

这样表述微调问题，可以为 LoRA 矩阵 $A$ 和 $B$ 找到更合适的量化初始化，也可以较好地近似原大语言模型的参数权重 $W$。为此，同时优化量化目标和 LoRA 微调目标。前者主要为了减少权重矩阵运算所需内存，后者主要为了在 $A$ 和 $B$ 的低秩约束下使任务表现尽可能好。

LoftQ 交替迭代，求解这个联合损失的最小值。先给定当前对 $A$ 和 $B$ 的估计，寻找使本步 Frobenius 范数最小的 $Q$。再给定新估计的 $Q$，用奇异值分解为 $A$ 和 $B$ 寻找量化权重 $Q$ 与原权重 $W$ 之间的残差 $Q-W$ 的低秩近似。通过交替估计量化值和执行量化感知的奇异值分解（singular value decomposition，SVD），LoftQ 在两者之间找到较好的折中，让它们共同改善微调结果。交替联合优化完成后，冻结最优的 $Q$，即可进行常规的 LoRA 微调。

LoftQ 在量化误差与 LoRA 低秩表示引入的误差之间取得了折中，而 QLoRA 并未明确最小化微调中的量化误差。由于量化相对原大语言模型引入精度偏差，QLoRA 的泛化表现不及 LoftQ。\refcite{ch04-li-2023a}{Li 等人（2023a）}测试的各项基准均显示 LoftQ 优于 QLoRA。

% PDF 191--197；原书第 162--168 页。
\section{终端用户提示}
\label{sec:ch04-end-user}
本章迄今讨论的学习策略，或是调优大语言模型的全部参数，或是调优部分参数，或是调优附加的适配器。它们都属于第~1.5.2~节介绍的大语言模型适配。接下来讨论终端用户提示（end-user prompting），即利用大语言模型的自回归文本生成和上下文学习能力，得到期望的结果（\refcite{ch04-minaee}{Minaee et al., 2024}；\refcite{ch04-zhao}{Zhao et al., 2023}）。

这些方法通常通过设计提示结构，在考虑大语言模型能力和限制的前提下，使输出尽可能满足具体应用的要求。设计提示，需要会理解和运用语言，尤其是相关领域的语言，还需要知道大语言模型的强项和弱项，并运用工程思维来组织和安排一条提示或一组提示中信息的顺序，从模型取得最有用的输出。与传统的数据科学和机器学习工程一样，提示工程既需要创造力，也需要仔细处理会影响结果的细节。

\begin{practicaltip}
可以把大语言模型的一次输出，想象成众多可能结果中的一个。提示是用户引导生成过程、使之偏向最有用结果的手段。在能力最强的大语言模型中，从单个子词词元到语法关系等更高层结构，用户都可以在这些层面影响模型输出，因为语言建模能够有效地学习这些结构（如 \refcite{ch04-jawahar}{Jawahar et al., 2019}）。有效的提示通常来自对内容和结构有步骤的试验。例如，考察领域专用词元或词语能否让大语言模型的回答更符合要求，或者比较正式语法和口语语法的影响（\refcite{ch04-chen}{Chen et al., 2023}）。
\end{practicaltip}

本章最后几节讨论常用的终端用户提示策略及其应用。我们介绍其中常见的技术和主要概念。其他技术可参阅 \refcite{ch04-chen}{Chen 等人（2023）}的综述及《提示工程指南》\footnote{\url{https://www.promptingguide.ai/}}。这些技术使用不同的结构安排，以更好地控制大语言模型输出与应用的契合程度。了解它们在不同条件下的效果，有助于在实际应用中用好大语言模型。

\subsection{零样本提示}
\label{sec:ch04-zero}
只包含任务指令的提示称为零样本提示。提示中没有额外的解题样例或示范，因此要仔细设计，才能激发完成目标任务所需的知识或能力。情感分类是一个例子，可采用清单~\ref{lst:ch04-zero} 的提示。

\noindent\begin{minipage}{\linewidth}
\begin{lstlisting}[caption={零样本情感分类提示},label={lst:ch04-zero}]
Please classify the following sentence as either 'Positive ',
'Neutral ' or 'Negative ' with respect to its sentiment.
Sentence: I hated the color of the front door!
Sentiment:
\end{lstlisting}
\end{minipage}\par

如前所述，零样本提示直接利用大语言模型已有的知识和能力。在清单~\ref{lst:ch04-zero} 的情感分类中，前提是大语言模型已经了解情感这一概念，以及文本是怎样表达情感的。

\subsection{少样本提示}
\label{sec:ch04-few}
如果零样本提示未能引出大语言模型已有的知识或能力，可以试用少样本提示。少样本提示既包含任务描述，也包含一个或多个解题示范。加入要模型完成的任务的示范，可以利用大语言模型的上下文学习能力，使任务表现优于零样本方法（\refcite{ch04-touvron}{Touvron et al., 2023}）。

仍以第~\ref{sec:ch04-zero}~节的情感分类任务为例，清单~\ref{lst:ch04-few} 给出一条少样本提示。

\noindent\begin{minipage}{\linewidth}
\begin{lstlisting}[caption={少样本情感分类提示},label={lst:ch04-few}]
Sentence: I just love it when I wake up to the sun shining
through my window.
Sentiment: Positive
Sentence: I was walking through the town yesterday.
Sentiment: Neutral
Sentence: I can't see a way to solve this problem without it
costing a lot.
Sentiment: Negative
Sentence: That sounds like such an exciting opportunity.
Sentiment:
\end{lstlisting}
\end{minipage}\par

\refcite{ch04-min}{Min 等人（2022）}发现，少样本提示的某些属性比较重要，其他属性的影响似乎较小。例如，清单~\ref{lst:ch04-few} 的提示采用固定的格式，反复出现 \texttt{Sentence} 后跟 \texttt{Sentiment} 的序列，以示范任务。对于任务表现，这个结构比示范本身是否正确更重要。即便给出的标签不正确，效果也可能优于完全不给标签。少样本提示对分类或实体抽取等任务很有效，但在需要复杂推理时有明显局限。接下来会介绍这类任务所用的思维链和思维树提示。

\subsection{提示链}
提示链（prompt chaining）把解决问题时与大语言模型的交互分成模块，使交互更简单。若单条提示不起作用，通常是因为问题太复杂，大语言模型难以用一条提示解决，这时就可以用提示链。把较大的问题拆成多条提示，再考虑它们的先后依赖，将其逐条连接起来，往往可以更好地控制结果，也能提高结果质量。

\noindent\begin{minipage}{\linewidth}
\begin{lstlisting}[caption={零样本情感分类提示},label={lst:ch04-finance-one}]
Please provide a short summary of the financial dealings
between each business entity pair within the following
document:
{{document}}
Summaries:
\end{lstlisting}
\end{minipage}\par

设想一项任务，需要为一篇文档中每对商业实体之间的各种财务往来写一段简短摘要。一种做法是构造清单~\ref{lst:ch04-finance-one} 的简单提示，让大语言模型在一次推理中解决整个问题。%
\footnote{译注：原书清单 4.5 至 4.7 的标题均为“零样本情感分类提示”，清单内容实际用于处理财务往来。}
在具体处理时，单条提示要求模型先理解指令，结合指令和文档作推理，再结合识别出的实体和文档作推理，最后为每对实体生成摘要。即使能力很强的大语言模型也可能难以完成。

\noindent\begin{minipage}{\linewidth}
\begin{lstlisting}[caption={零样本情感分类提示},label={lst:ch04-finance-pairs}]
Please list all business entity pairs within the following
document. Only entity pairs recorded in the document as
having had business dealings should be listed.
Document: {{document}}
Entity Pairs with business dealings:
\end{lstlisting}
\end{minipage}\par

考虑到任务的复杂性，可以采用提示链。先用清单~\ref{lst:ch04-finance-pairs} 的提示找出并列出文档中存在财务往来的所有商业实体对，再把结果交给后续的一条或多条提示，例如清单~\ref{lst:ch04-finance-summary} 的模板，用来分别概括各对实体的财务往来。这样做可能改善任务表现，也能更好地控制结果。把大问题拆成较小的任务模块后，开发人员可以评估大语言模型在中间步骤上的表现，只修改需要修改的步骤，就能改善整体表现。

\noindent\begin{minipage}{\linewidth}
\begin{lstlisting}[caption={零样本情感分类提示},label={lst:ch04-finance-summary}]
Please summarize the financial dealings between two entities
listed below, as recorded in the following document.
Entities: {{entity-pair}}
Document: {{document}}
Summary:
\end{lstlisting}
\end{minipage}\par

围绕提示链这一概念已经开发了多种框架，第~8~章会详细讨论。其中两种常用框架是 LangChain 和 DSPy，前者的抽象层次比后者高得多。它们旨在简化复杂提示链的开发，并使其开发流程更接近传统的软件开发实践。

\subsection{思维链}
\label{sec:ch04-cot}
\refcite{ch04-wei-2023}{Wei 等人（2023）}首先突出介绍了思维链（chain-of-thought，CoT）提示。它组织提示的上下文与示例，以模仿人解决问题时通常采用的逐步思考和推理过程。通常，能够自然分解为一系列中间问题的问题，很适合思维链提示。该方法中效果最好的提示，使用少量示例展示解决所提问题所需的推理步骤。思维链提示主要有三种变体。

\begin{itemize}
\item \textbf{零样本思维链。}\refcite{ch04-kojima}{Kojima 等人（2023）}提出，是三种变体中最简单、最直接的一种。只需在提示末尾加上“Let's think step by step”（我们逐步思考）或类似表述。\refcite{ch04-kojima}{Kojima 等人（2023）}发现，仅这一处改动，就能使大语言模型在 MultiArith 数据集（\refcite{ch04-roy}{Roy and Roth, 2015}）上的准确率从 17.7\% 提升到 78.7\%，在 GSM8K 数据集（\refcite{ch04-cobbe}{Cobbe et al., 2021}）上从 10.4\% 提升到 40.7\%。
\item \textbf{手工思维链。}提示由提示工程人员手工构造，含有一个或多个示范，表明希望大语言模型用哪些推理步骤解答样例。研究表明，在难度较大的数学问题上，加入这类示范可以达到当时的先进水平。
\item \textbf{自动思维链。}\refcite{ch04-zhang-2022}{Zhang 等人（2022）}提出的一种技术，可减少设计有效思维链提示所需的人工工作。示范越多样，思维链提示通常越有效，但手工构造这类提示可能很费力。自动思维链对问题聚类，再从各簇抽样，尽量增加示范问题的多样性，同时借助零样本思维链提示，让大语言模型生成这些示范的推理过程。将自动生成的示范纳入提示模板，再用模板推理。\refcite{ch04-zhang-2022}{Zhang 等人（2022）}表明，在相关基准上，这种方法的表现达到或超过了手工思维链提示。
\end{itemize}

\begin{practicaltip}
各种思维链提示策略都利用大语言模型涌现出的推理能力。研究表明，它们只有在模型超过一定规模（按参数数目计）时才有效。较小的大语言模型在任务表现上的提升不及更大的模型。例如，在 GSM8K 基准上，思维链提示相对于常规提示的最大提升，出现在具有 1750 亿参数的 GPT-3 模型上。常规提示的准确率为 15.6\%，思维链提示为 46.9\%。相较之下，具有 70 亿参数的 GPT-3 模型用这两种提示所得的准确率分别为 4\% 和 2.4\%（\refcite{ch04-wei-2023}{Wei et al., 2023}）。因此，开发人员需要验证思维链提示在所选大语言模型上是否有效。
\end{practicaltip}

\subsection{自一致性}
如前所述，大语言模型的输出容易出现编造或幻觉。在对一致性或事实准确性要求较高的应用中，自一致性提示（self-consistency prompting）是一种有效方法。一般的思路是，大语言模型对同一个查询给出的回答越一致，回答就越可能正确（\refcite{ch04-wang-2023a}{Wang et al., 2023a}）。

自一致性采用少样本思维链提示，多次以同一提示查询大语言模型，得到多个回答。再从这些回答中选出提示的答案，具体可以采用多种方法。在算术任务中，可用简单多数投票选择答案。对于问答等语言任务，可采用语义相似度或 $n$ 元组重叠等方法。

\subsection{思维树}
思维树（tree-of-thoughts，ToT）提示沿用思维链的基本思路，将大语言模型的注意力引向解题所需推理步骤的示范或描述。但思维树更接近人们寻找正确答案时所做的多路径探索（\refcite{ch04-long}{Long, 2023}）。它使大语言模型能沿着多条推理路径探索问题，摆脱线性推理链的限制。这样能减少因推导中的某一步出错而得到错误解的可能，也能通过探索更多解题路径来提高答对的可能。

\begin{keypoint}
思维树借鉴了关于人如何解题的研究。研究发现，人会在组合问题空间中进行认知搜索以找到解（\refcite{ch04-simon}{Simon and Newell, 1971}）。随着每个分支中的推理推进，一些路径被逐步排除，搜索范围越来越窄。思维树提示可以构造多条解题路径，还允许在这些路径上规划、前瞻和回溯，以确定效果最好的解题路线。
\end{keypoint}

思维树这一思想似乎由 \refcite{ch04-yao-2023}{Yao 等人（2023）}和 \refcite{ch04-long}{Long（2023）}独立提出。两项工作的主要区别，在于怎样在各个“思路”间搜索。前者采用广度优先搜索或深度优先搜索，后者采用通过强化学习训练的专用控制器模块。思维树可以看作自一致性的进一步扩展，除了可以通过多数投票选择答案，还可以抽样生成更多的中间推理步骤，最终得到正确答案。

% PDF 197--201；原书第 168--172 页。
\section{教程——在资源受限时微调大语言模型}
\label{sec:ch04-tutorial}
\subsection{概述}
前面已经介绍几种参数高效微调技术，并概述了大语言模型微调的两种主要方法，即指令调优和对齐调优。本教程用 LoRA 和 QLoRA 训练大语言模型，使它完成一项基于指令的特定任务。这里集中在一项任务上，覆盖范围比严格意义上的指令调优要窄，但构造模板时采用指令调优的方法。

\begin{keypoint}
\noindent\textbf{目标}
\begin{itemize}
\item 演示参数高效微调在内存需求和输出质量方面的优点。
\item 比较较大的大语言模型和规模较小的大语言模型的能力。
\item 为生成的文本制定评分标准，使用能力更强的大语言模型评分。
\end{itemize}
\end{keypoint}

这里给出教程的简明版本。完整版见 \url{https://github.com/springer-llms-deep-dive/llms-deep-dive-tutorials}。

\subsection{实验设计}
本教程要构建一个大语言模型，输入客户与客服人员的对话，输出概括主要内容的摘要。这里记录的结果来自一个配有 16~GB V100 GPU 的 Google Colab 会话。使用的数据集是 TWEETSUMM（\refcite{ch04-feigenblat}{Feigenblat et al., 2021}），由客户和不同公司客服人员在 \texttt{x.com}（原 Twitter）上来回交流的对话组成。每段对话配有人工撰写的两句摘要，记录客户的请求和客服人员的回复。大多数对话有不同标注者编写的多份摘要。

为评估大语言模型所生成摘要的质量，按下列三项标准计分。

\begin{enumerate}
\item 对客户问题或投诉的描述是否基本准确？
\item 对客服人员回复的描述是否基本准确？
\item 摘要是否恰好由两句话组成？
\end{enumerate}

每满足一项标准，摘要得一分。依照 \refcite{ch04-dettmers-2023}{Dettmers 等人（2023）}的做法，用 GPT-4 给摘要评分。把含有上述标准的评分细则连同输入对话和生成的摘要交给 GPT-4，让它返回一个满分为 3 分的分数。

首先测试 DistilGPT-2。这是一个有 8500 万参数的自回归大语言模型，由 GPT-2 提供监督来训练。选择它是因为内存需求相对较低，便于在本实验的 Colab 环境中微调。

随后尝试改用更大的大语言模型来改善结果。更丰富的语言知识有助于模型理解这些消息中发生了什么。为此，采用 Llama-2-7B，这是 Meta 于 2023 年发布的、具有 70 亿参数的自回归文本生成模型。虽然它的能力更强，但用微调 DistilGPT-2 的方法尝试微调它时，显存不足。这就需要采用参数高效微调技术。接下来使用 LoRA 和 QLoRA，比较不同训练方法得到的模型表现与训练时间。

\subsection{结果与分析}
\subsubsection{DistilGPT-2}
作为参照，先让 DistilGPT-2 在不做微调的情况下为测试集中的每段对话生成摘要。为文本生成定义一个 \texttt{transformers} 流水线，然后输入由 TWEETSUMM 测试集按模板构造的提示。输出质量较差。这并不意外，因为 DistilGPT-2 的规模较小，未经进一步微调，难以展现突出的涌现能力。接着用 Python 包 \texttt{trl} 在训练数据上微调模型。它对 \texttt{transformers} 的功能作了便于使用的封装。微调后的 DistilGPT-2 表现优于基础模型，尤其是在摘要长度这一项上，但对客户和客服人员对话的描述仍然不够准确。

为考察总体表现，分别用基础模型和微调模型为测试集中的 50 段对话生成摘要，再用 GPT-4 评分。基础模型摘要的累计得分为 2 分，满分为 150 分。这个结果很差，无法满足任务要求。微调模型的表现好得多，得到 67/150 分，不过仍与理想效果相差较远。

\begin{table}[!htbp]
\centering
\caption{在 TWEETSUMM 训练集上用各方法调优模型，并在测试集上推理得到的最终分数，满分为 150 分。表中同时列出各模型的调优时间。}
\label{tab:ch04-tweetsumm}
\small
\begin{tabularx}{\linewidth}{@{}Xrr@{}}
\toprule
模型配置 & 摘要得分（满分 150） & 调优时间（分钟）\\
\midrule
基础 DistilGPT-2 & 2 & 0\\
全参数微调的 DistilGPT-2 & 67 & 9.7\\
LoRA 调优的 DistilGPT-2 & 58 & 6.9\\
QLoRA 调优的 DistilGPT-2 & 52 & 14.3\\
基础 Llama-2-7B & 25.5 & 0\\
全参数微调的 Llama-2-7B & \multicolumn{2}{c@{}}{失败}\\
LoRA 调优的 Llama-2-7B & 131 & 75.1\\
QLoRA 调优的 Llama-2-7B & 125 & 21.3\\
\bottomrule
\end{tabularx}
\end{table}

如第~\ref{sec:ch04-lora}~节所述，低秩适配器是一种常用且高效的降低训练内存需求的方法。只需调优两个低秩矩阵，推理时将所得的权重更新加到原权重中，无需微调整个权重矩阵。这大幅减少了训练时需要在内存中保存梯度的参数数量。我们还试用效率更高的 QLoRA，先把模型权重量化到 4 位，再用 LoRA 方法调优。

LoRA 和 QLoRA 在 TWEETSUMM 数据集上的比较见表~\ref{tab:ch04-tweetsumm}。二者的表现都没有达到全参数微调的水平，但远好于基础模型。尽管 DistilGPT-2 的表现较低，这两种方法训练时的总 GPU 负担较小。与全参数微调相比，LoRA 调优的峰值 GPU 显存占用少 228~MB，QLoRA 调优少 336~MB。考虑到 DistilGPT-2 的权重矩阵约占 356~MB，这一差别相当明显。

\subsubsection{Llama-2-7B}
接着尝试用更大的语言模型来改善结果。它拥有更丰富的语言知识，可能更有能力理解这些消息中发生了什么。Llama-2-7B 符合要求。重复基础模型的零样本摘要实验，Llama-2 得到 25.5/150 分。表现依然较差，但较基础 DistilGPT-2 明显改善。接着尝试全参数微调 Llama-2，GPU 显存不足。Llama-2-7B 有 70 亿参数，仅模型权重便占约 12~GB 显存，加入微调所需的梯度后，总占用增至约 64~GB，远超本实验所用 V100 GPU 的 16~GB 显存。

\sourcefigure{fig04_12.pdf}{各种推理配置在 TWEETSUMM 摘要任务上的最终得分。上图显示原始分数，满分为 150 分。下图按三项标准拆分得分，即是否准确概括客户的发言、是否准确概括客服人员的发言，以及输出是否恰好两句话。由于显存不足，Llama-2-7B 的全参数微调未能产出模型。}{fig:ch04-tweetsumm}

因此，在 Llama-2-7B 上改用 PEFT 方法，训练得以开始，未出现 CUDA 错误。调优一轮，LoRA 用时 75 分钟，QLoRA 只用 21 分钟。这样做使表现大幅改善。测试集上，LoRA 调优模型得 131/150 分，QLoRA 调优模型得 125 分。%
\footnote{译注：原书正文将 QLoRA 的得分印为“125/131”。表~\ref{tab:ch04-tweetsumm} 和图~\ref{fig:ch04-tweetsumm} 均以 150 为满分，列出的 QLoRA 得分为 125。}

图~\ref{fig:ch04-tweetsumm} 汇总本教程各配置在测试集上的评估结果。两个经过适配器调优的 Llama-2-7B 模型总分最高，在每项评分标准上也表现最好。从下图可见，微调后的 DistilGPT-2 学会了将摘要限制为两句，但摘要仍不够准确，未达到 GPT-4 的评分要求。基础 Llama-2-7B 生成的摘要中，被判为准确的数量与全参数微调的 DistilGPT-2 相同，但没有经过进一步训练，就无法遵守格式要求。

\subsection{结论}
实验表明，较小的大语言模型可通过调优学会遵循特定指令，但受信息容量限制，其语义处理能力最终难与更大的模型相比。在两个经过调优的 Llama-2 模型中，QLoRA 略逊于 LoRA，但其调优时间不到后者的三分之一。对训练数据量较大的场景，这项权衡十分重要。总体而言，低秩适配器调优利用了 Llama-2-7B 的大量参数，得到了能够准确概括对话内容的摘要机器人。

% PDF 201--204；原书第 172--175 页。
\begin{chapterreferences}
\refentry{ch04-aghajanyan}{Armen Aghajanyan, Luke Zettlemoyer, and Sonal Gupta. Intrinsic dimensionality explains the effectiveness of language model fine-tuning, 2020.}
\refentry{ch04-brown}{Tom Brown, Benjamin Mann, Nick Ryder, Melanie Subbiah, Jared D Kaplan, Prafulla Dhariwal, Arvind Neelakantan, Pranav Shyam, Girish Sastry, Amanda Askell, et al. Language models are few-shot learners. \textit{Advances in Neural Information Processing Systems}, 33:1877–1901, 2020.}
\refentry{ch04-chen}{Banghao Chen, Zhaofeng Zhang, Nicolas Langrené, and Shengxin Zhu. Unleashing the potential of prompt engineering in large language models: A comprehensive review, 2023.}
\refentry{ch04-cobbe}{Karl Cobbe, Vineet Kosaraju, Mohammad Bavarian, Mark Chen, Heewoo Jun, Lukasz Kaiser, Matthias Plappert, Jerry Tworek, Jacob Hilton, Reiichiro Nakano, et al. Training verifiers to solve math word problems. \textit{arXiv preprint arXiv:2110.14168}, 2021.}
\refentry{ch04-dettmers-2022}{Tim Dettmers, Mike Lewis, Younes Belkada, and Luke Zettlemoyer. LLM.int8(): 8-bit matrix multiplication for Transformers at scale, 2022.}
\refentry{ch04-dettmers-2023}{Tim Dettmers, Artidoro Pagnoni, Ari Holtzman, and Luke Zettlemoyer. QLoRA: Efficient finetuning of quantized LLMs, 2023.}
\refentry{ch04-ding}{Ning Ding, Yujia Qin, Guang Yang, Fuchao Wei, Zonghan Yang, Yusheng Su, Shengding Hu, Yulin Chen, Chi-Min Chan, Weize Chen, et al. Parameter-efficient fine-tuning of large-scale pre-trained language models. \textit{Nature Machine Intelligence}, 5(3):220–235, 2023.}
\refentry{ch04-edalati}{Ali Edalati, Marzieh Tahaei, Ivan Kobyzev, Vahid Partovi Nia, James J. Clark, and Mehdi Rezagholizadeh. KronA: Parameter efficient tuning with Kronecker adapter, 2022.}
\refentry{ch04-feigenblat}{Guy Feigenblat, Chulaka Gunasekara, Benjamin Sznajder, Sachindra Joshi, David Konopnicki, and Ranit Aharonov. TWEETSUMM - a dialog summarization dataset for customer service. In \textit{Findings of the Association for Computational Linguistics: EMNLP 2021}, pages 245–260, Punta Cana, Dominican Republic, November 2021. Association for Computational Linguistics. URL \url{https://aclanthology.org/2021.findings-emnlp.24}.}
\refentry{ch04-frantar}{Elias Frantar, Saleh Ashkboos, Torsten Hoefler, and Dan Alistarh. GPTQ: Accurate post-training quantization for generative pre-trained Transformers, 2023.}
\refentry{ch04-he-2022}{Junxian He, Chunting Zhou, Xuezhe Ma, Taylor Berg-Kirkpatrick, and Graham Neubig. Towards a unified view of parameter-efficient transfer learning, 2022.}
\refentry{ch04-he-2015}{Kaiming He, Xiangyu Zhang, Shaoqing Ren, and Jian Sun. Delving deep into rectifiers: Surpassing human-level performance on ImageNet classification, 2015.}
\refentry{ch04-houlsby}{Neil Houlsby, Andrei Giurgiu, Stanislaw Jastrzebski, Bruna Morrone, Quentin de Laroussilhe, Andrea Gesmundo, Mona Attariyan, and Sylvain Gelly. Parameter-efficient transfer learning for NLP, 2019.}
\refentry{ch04-hu-2021}{Edward J. Hu, Yelong Shen, Phillip Wallis, Zeyuan Allen-Zhu, Yuanzhi Li, Shean Wang, Lu Wang, and Weizhu Chen. LoRA: Low-rank adaptation of large language models, 2021.}
\refentry{ch04-hu-2023}{Zhiqiang Hu, Lei Wang, Yihuai Lan, Wanyu Xu, Ee-Peng Lim, Lidong Bing, Xing Xu, Soujanya Poria, and Roy Ka-Wei Lee. LLM-Adapters: An adapter family for parameter-efficient fine-tuning of large language models, 2023.}
\refentry{ch04-huang}{Chengsong Huang, Qian Liu, Bill Yuchen Lin, Tianyu Pang, Chao Du, and Min Lin. LoRAHub: Efficient cross-task generalization via dynamic LoRA composition, 2023.}
\refentry{ch04-jawahar}{Ganesh Jawahar, Benoît Sagot, and Djamé Seddah. What does BERT learn about the structure of language? In \textit{Proceedings of the 57th Annual Meeting of the Association for Computational Linguistics}, pages 3651–3657, Florence, Italy, July 2019. Association for Computational Linguistics. doi: 10.18653/v1/P19-1356. URL \url{https://aclanthology.org/P19-1356}.}
\refentry{ch04-kim}{Sehoon Kim, Coleman Hooper, Amir Gholami, Zhen Dong, Xiuyu Li, Sheng Shen, Michael W. Mahoney, and Kurt Keutzer. SqueezeLLM: Dense-and-sparse quantization, 2023.}
\refentry{ch04-kingma}{Diederik P. Kingma and Jimmy Ba. Adam: A method for stochastic optimization, 2017.}
\refentry{ch04-kojima}{Takeshi Kojima, Shixiang Shane Gu, Machel Reid, Yutaka Matsuo, and Yusuke Iwasawa. Large language models are zero-shot reasoners, 2023.}
\refentry{ch04-kopiczko}{Dawid Jan Kopiczko, Tijmen Blankevoort, and Yuki Markus Asano. VeRA: Vector-based random matrix adaptation, 2023.}
\refentry{ch04-lester}{Brian Lester, Rami Al-Rfou, and Noah Constant. The power of scale for parameter-efficient prompt tuning. \textit{arXiv preprint arXiv:2104.08691}, 2021.}
\refentry{ch04-li-liang}{Xiang Lisa Li and Percy Liang. Prefix-tuning: Optimizing continuous prompts for generation. \textit{arXiv preprint arXiv:2101.00190}, 2021.}
\refentry{ch04-li-2023a}{Yixiao Li, Yifan Yu, Chen Liang, Pengcheng He, Nikos Karampatziakis, Weizhu Chen, and Tuo Zhao. LoftQ: LoRA-fine-tuning-aware quantization for large language models, 2023a.}
\refentry{ch04-li-2023b}{Yunxiang Li, Zihan Li, Kai Zhang, Ruilong Dan, Steve Jiang, and You Zhang. ChatDoctor: A medical chat model fine-tuned on a large language model Meta-AI (Llama) using medical domain knowledge, 2023b.}
\refentry{ch04-liu}{Zechun Liu, Barlas Oguz, Changsheng Zhao, Ernie Chang, Pierre Stock, Yashar Mehdad, Yangyang Shi, Raghuraman Krishnamoorthi, and Vikas Chandra. LLM-QAT: Data-free quantization aware training for large language models, 2023.}
\refentry{ch04-long}{Jieyi Long. Large language model guided tree-of-thought, 2023.}
\refentry{ch04-mahabadi}{Rabeeh Karimi Mahabadi, James Henderson, and Sebastian Ruder. Compacter: Efficient low-rank hypercomplex adapter layers, 2021.}
\refentry{ch04-min}{Sewon Min, Xinxi Lyu, Ari Holtzman, Mikel Artetxe, Mike Lewis, Hannaneh Hajishirzi, and Luke Zettlemoyer. Rethinking the role of demonstrations: What makes in-context learning work?, 2022.}
\refentry{ch04-minaee}{Shervin Minaee, Tomas Mikolov, Narjes Nikzad, Meysam Chenaghlu, Richard Socher, Xavier Amatriain, and Jianfeng Gao. Large language models: A survey, 2024.}
\refentry{ch04-ouyang}{Long Ouyang et al. Training language models to follow instructions with human feedback, 2022.}
\refentry{ch04-pfeiffer}{Jonas Pfeiffer, Aishwarya Kamath, Andreas Rücklé, Kyunghyun Cho, and Iryna Gurevych. AdapterFusion: Non-destructive task composition for transfer learning, 2021.}
\refentry{ch04-raffel}{Colin Raffel, Noam Shazeer, Adam Roberts, Katherine Lee, Sharan Narang, Michael Matena, Yanqi Zhou, Wei Li, and Peter J. Liu. Exploring the limits of transfer learning with a unified text-to-text Transformer, 2020.}
\refentry{ch04-roy}{Subhro Roy and Dan Roth. Solving general arithmetic word problems. In Lluís Màrquez, Chris Callison-Burch, and Jian Su, editors, \textit{Proceedings of the 2015 Conference on Empirical Methods in Natural Language Processing}, pages 1743–1752, Lisbon, Portugal, September 2015. Association for Computational Linguistics. doi: 10.18653/v1/D15-1202. URL \url{https://aclanthology.org/D15-1202}.}
\refentry{ch04-simon}{Herbert A Simon and Allen Newell. Human problem solving: The state of the theory in 1970. \textit{American Psychologist}, 26(2):145, 1971.}
\refentry{ch04-sung}{Yi-Lin Sung, Jaemin Cho, and Mohit Bansal. LST: Ladder side-tuning for parameter and memory efficient transfer learning, 2022.}
\refentry{ch04-taori}{Rohan Taori, Ishaan Gulrajani, Tianyi Zhang, Yann Dubois, Xuechen Li, Carlos Guestrin, Percy Liang, and Tatsunori B Hashimoto. Alpaca: A strong, replicable instruction-following model. Stanford Center for Research on Foundation Models. \url{https://crfm.stanford.edu/2023/03/13/alpaca.html}, 3(6):7, 2023.}
\refentry{ch04-thoppilan}{Romal Thoppilan, Daniel De Freitas, Jamie Hall, Noam Shazeer, Apoorv Kulshreshtha, Heng-Tze Cheng, Alicia Jin, Taylor Bos, Leslie Baker, Yu Du, et al. LaMDA: Language models for dialog applications. \textit{arXiv preprint arXiv:2201.08239}, 2022.}
\refentry{ch04-touvron}{Hugo Touvron et al. Llama 2: Open foundation and fine-tuned chat models, 2023.}
\refentry{ch04-wang-2023a}{Xuezhi Wang, Jason Wei, Dale Schuurmans, Quoc Le, Ed Chi, Sharan Narang, Aakanksha Chowdhery, and Denny Zhou. Self-consistency improves chain of thought reasoning in language models, 2023a.}
\refentry{ch04-wang-2023b}{Yizhong Wang, Yeganeh Kordi, Swaroop Mishra, Alisa Liu, Noah A. Smith, Daniel Khashabi, and Hannaneh Hajishirzi. Self-Instruct: Aligning language models with self-generated instructions, 2023b.}
\refentry{ch04-wei-2021}{Jason Wei, Maarten Bosma, Vincent Y Zhao, Kelvin Guu, Adams Wei Yu, Brian Lester, Nan Du, Andrew M Dai, and Quoc V Le. Finetuned language models are zero-shot learners. \textit{arXiv preprint arXiv:2109.01652}, 2021.}
\refentry{ch04-wei-2023}{Jason Wei, Xuezhi Wang, Dale Schuurmans, Maarten Bosma, Brian Ichter, Fei Xia, Ed Chi, Quoc Le, and Denny Zhou. Chain-of-thought prompting elicits reasoning in large language models, 2023.}
\refentry{ch04-wortsman}{Mitchell Wortsman, Tim Dettmers, Luke Zettlemoyer, Ari Morcos, Ali Farhadi, and Ludwig Schmidt. Stable and low-precision training for large-scale vision-language models, 2023.}
\refentry{ch04-yao-2023}{Shunyu Yao, Dian Yu, Jeffrey Zhao, Izhak Shafran, Thomas L. Griffiths, Yuan Cao, and Karthik Narasimhan. Tree of thoughts: Deliberate problem solving with large language models, 2023.}
\refentry{ch04-yao-2022}{Zhewei Yao, Reza Yazdani Aminabadi, Minjia Zhang, Xiaoxia Wu, Conglong Li, and Yuxiong He. ZeroQuant: Efficient and affordable post-training quantization for large-scale Transformers, 2022.}
\refentry{ch04-zaken}{Elad Ben Zaken, Shauli Ravfogel, and Yoav Goldberg. BitFit: Simple parameter-efficient fine-tuning for Transformer-based masked language-models, 2022.}
\refentry{ch04-zhang-2023a}{Boyu Zhang, Hongyang Yang, and Xiao-Yang Liu. Instruct-FinGPT: Financial sentiment analysis by instruction tuning of general-purpose large language models, 2023a.}
\refentry{ch04-zhang-2023b}{Qingru Zhang, Minshuo Chen, Alexander Bukharin, Pengcheng He, Yu Cheng, Weizhu Chen, and Tuo Zhao. Adaptive budget allocation for parameter-efficient fine-tuning, 2023b.}
\refentry{ch04-zhang-2023c}{Shengyu Zhang et al. Instruction tuning for large language models: A survey, 2023c.}
\refentry{ch04-zhang-2023d}{Yue Zhang, Leyang Cui, Deng Cai, Xinting Huang, Tao Fang, and Wei Bi. Multi-task instruction tuning of Llama for specific scenarios: A preliminary study on writing assistance, 2023d.}
\refentry{ch04-zhang-2022}{Zhuosheng Zhang, Aston Zhang, Mu Li, and Alex Smola. Automatic chain of thought prompting in large language models, 2022.}
\refentry{ch04-zhao}{Wayne Xin Zhao et al. A survey of large language models, 2023.}
\refentry{ch04-zhu}{Yaoming Zhu, Jiangtao Feng, Chengqi Zhao, Mingxuan Wang, and Lei Li. Counter-interference adapter for multilingual machine translation, 2021.}
\end{chapterreferences}

